\documentclass[../main.tex]{subfiles}

\begin{document}

\onlyinsubfile{
    \renewcommand{\contentsname}{ОГЛАВЛЕНИЕ}
    \setcounter{tocdepth}{3}
    \tableofcontents
}

\refstepcounter{chapter}
\chapter*{\MakeUppercase{Глава \arabic{chapter}\\ Теоретические основы и базовые методы генерации молекул}}\label{chapter:base}
\addcontentsline{toc}{chapter}{Глава \arabic{chapter} Теоретические основы и базовые методы генерации молекул}

Разработка новых лекарственных препаратов с помощью методов глубокого обучения требует решения двух фундаментальных задач: представления сложной молекулярной структуры в формате, понятном для нейронной сети, и создания генеративных моделей, способных оперировать этим представлением для создания новых, химически осмысленных соединений. Данная глава посвящена рассмотрению теоретических основ, лежащих в основе этих процессов. Будет проанализирована линейная нотация SMILES, ставшая стандартом де-факто для представления молекул в виде последовательностей, а также рассмотрены базовые архитектуры, такие как рекуррентные нейронные сети (RNN) и вариационные автоэнкодеры (VAE), которые первыми продемонстрировали свою эффективность в задачах \textit{de novo} дизайна молекул.

% Содержимое подраздела 1.1 - Представление молекулярных структур для задач машинного обучения

\section{Представление молекулярных структур для задач машинного обучения}
\label{sec:representation}

Ключевым шагом, позволившим применить методы обработки естественного языка (Natural Language Processing, NLP) к задачам медицинской химии, стала разработка линейных нотаций для представления молекулярных структур. Наиболее широкое распространение получила система \textbf{SMILES (Simplified Molecular-Input Line-Entry System)} --- формат, представляющий молекулу в виде строки символов ASCII. Данный подход был разработан с целью создания машиночитаемого, но при этом компактного и интуитивно понятного для человека представления~\cite{wigh2022review}.

Принцип кодирования структуры в SMILES основан на наборе простых правил:
\begin{itemize}
	\item \textbf{Атомы} представляются стандартными символами из периодической таблицы (например, \texttt{C} для углерода, \texttt{N} для азота). Атомы в ароматических системах обозначаются строчными буквами (например, бензол кодируется как \texttt{c1ccccc1}).
	\item \textbf{Связи} одинарные подразумеваются по умолчанию. Двойные и тройные связи обозначаются символами \texttt{=} и \texttt{\#} соответственно.
	\item \textbf{Ветвления} боковых цепей от основного скелета молекулы заключаются в круглые скобки \texttt{()}. Например, изобутан представляется строкой \texttt{CC(C)C}.
	\item \textbf{Циклы} кодируются путем разрыва одной из связей в цикле. Атомы, соединенные этой связью, помечаются одинаковыми цифровыми индексами. Например, циклогексан можно представить как \texttt{C1CCCCC1}, где цифра \texttt{1} указывает на замыкание цикла между первым и последним атомами углерода.
\end{itemize}

Представление SMILES обладает рядом преимуществ, главным из которых является его совместимость с моделями глубокого обучения. Рассматривая молекулу как <<предложение>> на <<химическом языке>>, исследователи смогли напрямую применить к задачам дизайна лекарств мощные архитектуры, такие как рекуррентные нейронные сети (RNN) и трансформеры. Кроме того, SMILES-строки компактны и более читаемы для человека по сравнению с другими форматами, такими как таблицы связности (connection tables), что упрощает работу с базами данных~\cite{wigh2022review}.

Тем не менее, у данного подхода есть и существенные недостатки. Основной проблемой является \textbf{неканоничность}: одна и та же молекула может быть представлена множеством синтаксически корректных SMILES-строк (например, \texttt{CCO} и \texttt{OCC} для этанола). Это усложняет сравнение структур и требует вычислительно затратных процедур канонизации. Другой важный недостаток заключается в том, что генеративные модели, обученные на SMILES, часто производят синтаксически некорректные строки (например, с непарными скобками). Это создает риск того, что модель в первую очередь изучает синтаксис нотации, а не фундаментальные химические закономерности. Наконец, в базовой форме SMILES не кодирует 3D-информацию о конформации молекулы, что является критически важным для многих задач.

Для решения этих проблем были разработаны альтернативные подходы. \textbf{InChI (International Chemical Identifier)} генерирует единственную, каноническую строку для любой молекулы, что решает проблему неоднозначности, но его сложная иерархическая структура делает его менее удобным для генеративных моделей. Более фундаментальной альтернативой являются \textbf{молекулярные графы}, где атомы представляются узлами, а связи -- ребрами. Такое представление является более естественным для молекул, лишено синтаксических артефактов и по построению не может сгенерировать химически невалидную структуру (например, с <<непарной связью>>). Однако работа с графовыми данными требует применения специализированных архитектур, таких как графовые нейронные сети.

% Содержимое подраздела 1.2 - Рекуррентные нейронные сети как генеративные модели "химического языка"


\section{Моделирование молекулярных структур как последовательностей с использованием рекуррентентных нейронных сетей}
\label{sec:rnn}

Исторически первыми и наиболее влиятельными архитектурами для генерации молекул стали рекуррентные нейронные сети (RNN), в частности их более сложная разновидность -- сети с долгой краткосрочной памятью (Long Short-Term Memory, LSTM). Принцип их работы в контексте химии основан на концепции авторегрессионной языковой модели. SMILES-строка разбивается на последовательность токенов (отдельных символов), и модель обучается на каждом шаге предсказывать следующий токен на основе всех предыдущих. В процессе генерации модель, начиная со стартового токена, посимвольно <<пишет>> новую SMILES-строку, постоянно опираясь на уже сгенерированную часть последовательности~\cite{wigh2022review}. Ранние работы в этой области продемонстрировали высокую эффективность, масштабируемость и, что самое главное, практическую применимость данного подхода.

Одним из первых масштабных исследований, доказавших состоятельность RNN-подхода, стала работа Ertl et al. (2017)~\cite{ertl2017silico}. Обучив относительно простую двухслойную LSTM-сеть на 509,000 биоактивных молекул из базы данных ChEMBL, авторы продемонстрировали возможность генерации \textbf{одного миллиона химически корректных молекул менее чем за два часа}. Несмотря на то, что только 32\% от всех сгенерированных строк оказались валидными, новизна результата была чрезвычайно высока: лишь 0.3\% сгенерированных структур совпадали с молекулами из обучающего набора. Ключевым достижением работы стала многоуровневая валидация, которая показала, что распределения физико-химических свойств (молекулярный вес, LogP, TPSA) и оценки синтетической доступности (SA score) у сгенерированных молекул практически идентичны реальным соединениям из ChEMBL. Наиболее убедительным доказательством стало применение QSAR-моделей, которое показало, что новые молекулы имеют сопоставимый с известными лекарствами потенциал биологической активности.

В то время как работа Ertl et al. продемонстрировала масштабируемость, исследование Bjerrum и Threllfall (2017)~\cite{bjerrum2017molecular} было сфокусировано на детальном анализе качества и новизны генерируемых структур. Обучая модель на различных подмножествах базы ZINC (<<фрагменто-подобные>> и <<лекарственно-подобные>> молекулы), они показали, что RNN способна точно воспроизводить распределения свойств из обучающей выборки. Это доказывает, что модель улавливает скрытые <<химические правила>>, заложенные в данных. Анализ новизны показал, что до \textbf{83\%} сгенерированных молекул отсутствовали в исходном обучающем наборе, при этом их оценка синтетической доступности (SA score) имела такое же распределение, как и у реальных соединений. Таким образом, было доказано, что модель генерирует не просто валидные, но и химически осмысленные, новые и потенциально синтезируемые структуры.

Кульминацией ранних исследований стала работа, продемонстрировавшая переход от теоретических выкладок к реальной экспериментальной проверке~\cite{merk2018novo}. В этом исследовании был применен подход на основе переноса обучения (transfer learning): RNN-модель, предварительно обученная на большой базе данных ChEMBL, была дообучена на очень малом наборе из 25 молекул-агонистов ядерных рецепторов RXR и PPAR. Из сгенерированных кандидатов пять были отобраны для химического синтеза. В результате \textbf{четыре из пяти соединений продемонстрировали желаемую биологическую активность} в наномолярном и микромолярном диапазоне, причем одно из соединений показало активность \texttt{EC\textsubscript{50} = 0.06 $\mu$M}. Эта работа стала первым проспективным применением генеративной модели, завершившим полный цикл от дизайна с помощью ИИ до синтеза и успешной биологической валидации, и доказала практическую ценность RNN-подхода для de novo дизайна лекарств.

% Содержимое подраздела 1.3 - Генеративные модели на основе вариационных автоэнкодеров


\section{Генеративные модели на основе вариационных автоэнкодеров}
\label{sec:vae}

Несмотря на успехи RNN, они обладают существенным недостатком: процесс генерации сложно контролировать. Отсутствует явный механизм, позволяющий целенаправленно создавать молекулы с заданными свойствами. Для решения этой проблемы были предложены модели на основе \textbf{вариационных автоэнкодеров (Variational Autoencoder, VAE)}. VAE состоит из двух частей: энкодера, который сжимает дискретное представление молекулы (например, SMILES) в непрерывный вектор в скрытом (латентном) пространстве, и декодера, который восстанавливает из этого вектора исходную молекулу. Главное преимущество VAE заключается в том, что его латентное пространство является непрерывным и структурированным, что открывает возможность для градиентной оптимизации свойств и интерполяции между известными структурами для создания новых~\cite{wigh2022review}.

Однако стандартные VAE сталкиваются с рядом проблем, главной из которых является \textbf{коллапс апостериорного распределения (posterior collapse)}. Этот эффект возникает, когда декодер учится игнорировать информацию из латентного пространства и генерирует выходные данные, опираясь только на входные условия, что приводит к потере разнообразия. Для решения этой проблемы в работе, представляющей архитектуру \textbf{PCF-VAE}~\cite{bhadwal2025pcf}, была предложена модифицированная функция потерь. В ней вводится динамический весовой коэффициент $\alpha$ для члена KL-дивергенции, который постепенно увеличивается в процессе обучения (метод <<отжига>>). Это позволяет модели на ранних этапах сфокусироваться на точности реконструкции, а затем плавно <<включать>> регуляризацию латентного пространства. В сочетании с использованием более робастной нотации GenSMILES, данный подход позволил достичь на бенчмарке MOSES высокой валидности (98.01\%), уникальности (100\%) и новизны (93.77\%), демонстрируя эффективное решение проблемы коллапса.

Другой важной задачей является создание моделей для условной генерации, то есть создания молекул с заранее определенными свойствами. В работе, описывающей \textbf{условный $\beta$-VAE}~\cite{richards2022conditional}, предложен элегантный метод <<явного обусловливания латентного пространства>> (Explicit Latent Conditioning, ELC). Вместо того чтобы центрировать априорное распределение в латентном пространстве вокруг нуля, модель центрирует его вокруг стандартизированного значения целевого свойства $\hat{c}$, то есть использует распределение $N(\hat{c}, I)$. Это заставляет модель организовывать латентное пространство в соответствии с целевым свойством, группируя молекулы со схожими значениями вместе. Данный подход, реализованный в гибридной архитектуре (Transformer-энкодер и GRU-декодер), показал рекордные результаты в задаче неограниченной оптимизации, сгенерировав молекулу с предсказанным значением penalized LogP (pLogP) равным \textbf{104.29}, что на порядок превосходило предыдущие достижения.

Несмотря на успехи строковых моделей, проблема генерации невалидных SMILES остается актуальной. Это привело к развитию графовых VAE, которые работают не с линейной строкой, а с более естественным графовым представлением молекулы. Одним из наиболее успешных примеров является \textbf{Junction Tree VAE (JT-VAE)}~\cite{kondratyev2023generative}. Эта модель оперирует на уровне химически осмысленных строительных блоков (колец и связей), что по построению гарантирует \textbf{100\% химическую валидность} генерируемых структур. Такой подход позволяет избежать стадии валидации и сосредоточиться на оптимизации свойств в хорошо структурированном латентном пространстве.


\onlyinsubfile{
    \printbibliography[heading=bibintoc, title={Библиографический список}]
}

\end{document}